\documentclass[12pt]{article}
\usepackage[utf8]{inputenc}
\usepackage{amsmath,amssymb}
\usepackage{graphics}
\usepackage{pbox}
\usepackage{tikz, tkz-fct, pgfplots}
\usetikzlibrary{arrows}
\usepackage{geometry}
\geometry{
	a4paper,
	total={170mm,257mm},
	left=20mm,
	top=20mm
} 
\usepackage[labelsep=period]{caption}


% ----------------- Commands -----------------
\renewcommand{\baselinestretch}{1.3}
\newcommand{\eps}{\varepsilon}
\newcommand\tline[2]{$\underset{\text{#1}}{\text{\underline{\hspace{#2}}}}$}
\DeclareMathOperator{\arccot}{arccot}


% ----------------- Main Document -----------------
\begin{document}
\pagestyle{empty}
\centerline{\large Ministry of Education and Science}	
\centerline{\large Federal State Budgetary Educational}
\centerline{\large Institution of Higher Education}
\centerline{\large ``Bauman Moscow State Technical University''}
\centerline{\large (BMSTU)}
\vspace{0.2cm}
\hrule
\vspace{0.5cm}
\begin{figure}[h]
\center
\includegraphics[height=0.35\linewidth]{img/bmstu-logo-color.png}
\end{figure}
\begin{center}
	\large	
	\begin{tabular}{c}
		Faculty of ``Fundamental Sciences'' \\
		Department of ``Higher Mathematics''		
	\end{tabular}
\end{center}
\vspace{0.5cm}
\begin{center}
	\LARGE \bf	
	\begin{tabular}{c}
		\textsc{Coursework Report} \\
		1-st semester  of the \\
		2020---2021 academic year.
	\end{tabular}
\end{center}
\vspace{0.5cm}
\begin{center}
	\large
	\begin{tabular}{p{6cm}ll}
		\pbox{6.1cm}{
			Coursework supervisor,\\
			sn. professor of the FS1 dep.} 	& \tline{\it(signature)}{5cm} & O. V. Kravchenko \\[0.8cm]
		Student of the FS1-11B\ \ academic group 		& \tline{\it(signature)}{5cm} & M. M. Alekseev
	\end{tabular}
\end{center}
\vfill
\begin{center}
	\large	
	\begin{tabular}{c}
		Moscow, \\
		2020.
	\end{tabular}
\end{center}


\newpage	
\tableofcontents

\newpage
\pagestyle{plain}
\section{Goals and objectives of the coursework}	
\subsection{Goals}
--- development of capabilities necessary for future professional activity and auxiliary for the successful completion of the baccalaureate course.

\subsection{Objectives}
\begin{enumerate}
\item Development of familiarity with the software necessary for future professional activity.
\item Development of the ability to research required information in specialised literature and other sources.
\item Development of the skills needed to formulate reports and present results.
\end{enumerate}

\subsection{Individual objectives}	
\begin{enumerate}
\item Examine the methods used to display mathematical information in the \LaTeX {} typesetting system.  
\item Examine the possibilities and potential of the Git version control system.
\item Learn to typeset texts with mathematical context, containing formulas and graphs, in the \LaTeX {} system. For this, complete the installation of the open-source distributive TeXLive and compiler TeXStudio.
\item Recreate in the \LaTeX {} system your variant of the sample calculations from the Mathematical Analysis course.
\item Register on the online platform GitHub and upload the initial tex-files and the results of their compilation (in pdf format).  
\end{enumerate} 

\newpage
\section{Report}
The relevance of this topic is dictated by the fundamental necessity of the ability to operate the typesetting systems of \LaTeX {} and TexStudio for the ultimate scope of displaying text, mathematical formulae and graphs. The skills acquired through the completion of this coursework may be applied in such processes as the writing of course projects, thesis work and future professional activity.

Moreover, the \LaTeX {} typesetting system allows for the usage of a large number of external instrumental packages, which greatly facilitate the process of displaying information in various fields of engineering, technical and scientific activity.
\newpage
\section{Individual objectives}

\subsection{Limits and Continuity.}
% ---------------------------- Task 1 ----------------------------------
\subsection*{\center Question \#1.}
{\bf Conditions:~}
Consider the sequence $a_n=\dfrac{7n-1}{n+1}$ and the value $c=7$. 
Prove that $\lim\limits_{n\rightarrow\infty}a_n=c$. 
More precisely, $\forall\eps>0$ find the minimum value $N=N(\eps)$ such that $\left|a_n-c\right|<\eps,\ \forall n>N(\eps)$.
Fill in the following table.
\begin{center}
	\begin{tabular}{|c|c|c|c|}
		\hline
		$\eps$ &  $0{,}1$ & $0{,}01$ & $0{,}001$ \\
		\hline
		$N(\eps)$ & & & \\
		\hline
	\end{tabular}
\end{center}
{\bf Solution:~}
If $c$ is the limit of $a_n$, then $\forall\eps>0$ and $\forall n>N(\eps)$, $\left|a_n-c\right|<\eps$ must hold true.
By substituting the values for $a_n$ and $c$ into the above inequality we achieve the following:
$$\left|\dfrac{7n-1}{n+1}-7\right|<\eps.$$
Next we open the modulus to form a double inequality and rewrite the difference to have a common denominator.
$$-\eps<\dfrac{-8}{n+1}<\eps.$$
Obviously, the right inequality holds true $\forall n \in \mathbb{N}$, therefore form this point on we will only consider the left inequality:
$$-\eps<\dfrac{-8}{n+1}.$$
At this point, by performing a series of transformations we can isolate $n$ and formulate a formula for $N(\eps)$.
$$
\begin{array}{c}
-\eps<\dfrac{-8}{n+1},\\ [12pt]
\eps>\dfrac{8}{n+1},\\ [12pt]
n+1>\dfrac{8}{\eps},\\ [12pt]
n>\dfrac{8}{\eps}-1,\\ [12pt]
N(\eps)=\left\lfloor\dfrac{8}{\eps}-1\right\rfloor.
\end{array}
$$
Where $\lfloor\ \ \rfloor$ {} is the floor function.\\
The existence of a $N=N(\eps)$ function proves that $\left|a_n-c\right|<\eps$ holds true for $\forall\eps>0$ and $\forall n>N(\eps) \iff \lim\limits_{n\rightarrow\infty}a_n=c$. We may fill in the table:
\begin{center}
	\begin{tabular}{|c|c|c|c|}
		\hline
		$\eps$ &  $0{,}1$ & $0{,}01$ & $0{,}001$ \\
		\hline
		$N(\eps)$ & 79 & 799  &  7999 \\
		\hline
	\end{tabular}
\end{center}
{\bf Verification:}
$$
\begin{array}{l}
\left|a_{80}-c\right|=\dfrac{-8}{81}<0.1,  \\[12pt]
\left|a_{800}-c\right|=\dfrac{-8}{801}<0.01,  \\[12pt]
\left|a_{8000}-c\right|=\dfrac{-8}{8001}<0.001.
\end{array}
$$

% ---------------------------- Task 2 ----------------------------------
\subsection*{\center Question \#2.}
{\bf Conditions:~} 
Evaluate the limits of the following functions:
$$
\begin{array}{cc}
\text{\bf(a)} & \lim\limits_{x\rightarrow-1}\dfrac{x^3-3x-2}{x^2-x-2}, \\[12pt]
\text{\bf(b)} & \lim\limits_{x\rightarrow+\infty}\dfrac{1-\sqrt[3]{4x^4-x^7\sqrt{x}}}{2x^2-3x+5}, \\[12pt]
\text{\bf(c)} & \lim\limits_{x\rightarrow3}\dfrac{\sqrt{x+13}-2\sqrt{x+1}}{x^2-9}, \\[12pt]
\text{\bf(d)} & \lim\limits_{x\rightarrow0}({2-\cos{3x}})^{\frac{1}{\ln{(1+x^2)}}}, \\[12pt]
\text{\bf(e)} & \lim\limits_{x\rightarrow0}{\left(\dfrac{\sin{2x}}{\arcsin{3x}}\right)}^{\arccot x}, \\[12pt]
\text{\bf(f)} & \lim\limits_{x\rightarrow2}\dfrac{\lg{(5-2x)}}{\sqrt{10-3x}-2}.
\end{array}
$$
{\bf Solutions:}\\
{\bf\flushleft (a)}
$$
\begin{array}{l}
\lim\limits_{x\rightarrow-1}\dfrac{x^3-3x-2}{x^2-x-2}= 
\left[\dfrac{0}{0}\right] = 
\lim\limits_{x\rightarrow-1}\dfrac{(x+1)(x^2-x-2)}{x^2-x-2} = \lim\limits_{x\rightarrow-1}(x+1) = 0.
\end{array}
$$
\vspace{0.2cm}
{\bf\flushleft (b)}
$$
\begin{array}{l}
\lim\limits_{x\rightarrow+\infty}\dfrac{1-\sqrt[3]{4x^4-x^7\sqrt{x}}}{2x^2-3x+5} =
\left[\dfrac{1-\sqrt[3]{\infty-\infty}}{\infty-\infty}\right] =
\lim\limits_{x\rightarrow+\infty}\dfrac{1-x^{5/2}\cdot\sqrt[3]{4x^{-7/2}-1}}{x^2\cdot(2-3x^{-1}+5x^{-2})} = \\ [15pt]
\lim\limits_{x\rightarrow+\infty}\dfrac{x^{5/2}\cdot(x^{-5/2}-\sqrt[3]{4x^{-7/2}-1})}{x^2\cdot(2-3x^{-1}+5x^{-2})} =
\lim\limits_{x\rightarrow+\infty}\sqrt{x}\cdot\dfrac{x^{-5/2}-\sqrt[3]{4x^{-7/2}-1}}{2-3x^{-1}+5x^{-2}} = +\infty.
\end{array}
$$
\vspace{0.2cm}
{\bf\flushleft (c)}
$$
\begin{array}{l}
\lim\limits_{x\rightarrow3}\dfrac{\sqrt{x+13}-2\sqrt{x+1}}{x^2-9} =
\left[\dfrac{0}{0}\right] =
\lim\limits_{x\rightarrow3}\dfrac{(x+13)-4(x+1)}{(x^2-9)(\sqrt{x+13}+2\sqrt{x+1})} = \\ [15pt]
\lim\limits_{x\rightarrow3}\dfrac{-3x+9}{(x-3)(x+3)(\sqrt{x+13}+2\sqrt{x+1})} =
\lim\limits_{x\rightarrow3}\dfrac{-3}{(x+3)(\sqrt{x+13}+2\sqrt{x+1})} = -\dfrac{1}{16}.
\end{array}
$$
\newpage
{\bf\flushleft (d)}
$$
\begin{array}{l}
\lim\limits_{x\rightarrow0}({2-\cos{3x}})^{\frac{1}{\ln{(1+x^2)}}} =
\left[1^{\infty}\right] =
\lim\limits_{x\rightarrow0}\exp\left(\dfrac{\ln(2-\cos{3x})}{\ln{(1+x^2)}}\right) =
\left|\begin{array}{l}
\ln(2-\cos{3x})\sim 1-\cos{3x} \\
\ln(1+x^2)\sim x^2
\end{array}\right| = \\[16pt]
\exp\left(\lim\limits_{x\rightarrow0}\dfrac{1-\cos{3x}}{x^2}\right) = 
\left|\begin{array}{l}
1-\cos{3x}\sim \dfrac{(3x)^2}{2}
\end{array}\right| = 
\exp\left(\lim\limits_{x\rightarrow0}\dfrac{9x^2}{2x^2}\right) =
e^{\frac{9}{2}} = (\sqrt{e})^9.
\end{array}
$$
\vspace{0.2cm}
{\bf\flushleft (e)}
$$
\begin{array}{l}
\lim\limits_{x\rightarrow0}{\left(\dfrac{\sin{2x}}{\arcsin{3x}}\right)}^{\arccot x} =
\left[\left(\dfrac{0}{0}\right)^{\frac{\pi}{2}}\right] =
\left|\begin{array}{l}
\sin{2x}\sim 2x \\
\arcsin{3x}\sim 3x
\end{array}\right| =
\left(\dfrac{2}{3}\right)^{\lim\limits_{x\rightarrow0}\arccot x} =
\left(\dfrac{2}{3}\right)^{\pi/2}.
\end{array}
$$
\vspace{0.2cm}
{\bf\flushleft (f)}
$$
\begin{array}{l}
\lim\limits_{x\rightarrow2}\dfrac{\lg{(5-2x)}}{\sqrt{10-3x}-2} =
\left[\dfrac{0}{0}\right] =
\left|\begin{array}{l}
t=x-2 \\
t\rightarrow0
\end{array}\right| =
\lim\limits_{t\rightarrow0}\dfrac{\lg{(1-2t)}}{\sqrt{4-3t}-2} =
\lim\limits_{t\rightarrow0}\ \dfrac{\lg{(1-2t)}}{2\cdot\sqrt{1-\frac{3t}{4}}-2} = \\ [15pt]
\left|\begin{array}{l}
\lg{(1-2t)}\sim \dfrac{-2t}{\ln{10}} \\ [10pt]
\sqrt{1-\frac{3t}{4}}-1\sim \dfrac{1}{2}\cdot\dfrac{-3t}{4}
\end{array}\right| =
\dfrac{1}{2\ln{10}}\cdot\lim\limits_{t\rightarrow0}\dfrac{16t}{3t} = \dfrac{8}{3\ln{10}}.
\end{array}
$$

% ---------------------------- Task 3 ----------------------------------
\subsection*{\center Question \#3.}
{\bf Conditions:}\\
(a) Show that the given functions $f(x)=x^2+\sqrt{x}+\sin{x}$ and $g(x)=\ln{\cos{\sqrt{x}}}$ are infinitely-large or infinitely-small functions as $x\rightarrow0+$;\\
(b) For both functions find their main part (equivalent function in the form $C(x-x_0)^\alpha$ if $x\rightarrow x_0$ or $C(x)^\alpha$ if $x\rightarrow\infty$), state their asymptotic orders of growth;\\
(c) Compare the functions $f(x)$ and $g(x)$ as $x\rightarrow0+$.\\ [0.5cm]
{\bf Solutions:}
{\bf\flushleft (a)}
To show that the given are functions are infinitely-small as $x\rightarrow0+$ we should evaluate their limits with the given tendency of the argument $x$:
$$
\begin{array}{c}
\lim\limits_{x\rightarrow0+}f(x) = 
\lim\limits_{x\rightarrow0+}(x^2+\sqrt{x}+\sin{x}) =
\left[0+0+0\right] = 0. \\ [13pt]
\lim\limits_{x\rightarrow0+}g(x) = 
\lim\limits_{x\rightarrow0+}\ln{\cos{\sqrt{x}}} =
\left[\ln{\cos{0}}\right] = 
\left[\ln{1}\right] = 0.
\end{array}
$$
\newpage
{\bf\flushleft (b)}
The functions $f(x)$ and $g(x)$ are infinitely-small as $x\rightarrow0+$, therefore their main parts have the form of $C(x)^\alpha$, and may be found using asymptotically equivalent functions.
$$
\begin{array}{c}
\sin{x}\sim x,\ \text{as}\ x\rightarrow{0+} \implies
f(x)=x^2+\sqrt{x}+\sin{x}\ \sim\ x^2+\sqrt{x}+x, \\
\because \text{ an infinitely-small polynomial is equivalent to its monomial of the lowest power:} \\
x^2+\sqrt{x}+x\ \sim\ \sqrt{x}=x^{(1/2)}, \\
\therefore\ f(x)=x^2+\sqrt{x}+\sin{x}\ \sim\ x^{(1/2)}.\\ [0.4cm]
1-\cos{x}\sim \dfrac{x^2}{2},\ \text{as}\ x\rightarrow{0+} \implies
g(x)=\ln{\cos{\sqrt{x}}}=\ln{(1+(\cos{\sqrt{x}}-1))}\ \sim\ \ln{(1-\dfrac{x}{2})},\\
\ln{(1+x)}\sim x,\ \text{as}\ x\rightarrow{0+} \implies
\ln{(1-\dfrac{x}{2})}\sim -\dfrac{x}{2}, \\
\therefore\ g(x)=\ln{\cos{\sqrt{x}}}\ \sim\ -\dfrac{x}{2}.
\end{array}
$$
Consequently, as $x\rightarrow0+$, the main part of $f(x)$ is $1(x)^{1/2}$ and the main part of $g(x)$ is $-\dfrac{1}{2}(x)^1$.
{\bf\flushleft (c)} To compare $f(x)$ and $g(x)$ as $x\rightarrow0+$, we should consider the limit of the relation between the main parts of $f(x)$ and $g(x)$:
$$
\lim\limits_{x\rightarrow0+}\dfrac{g(x)}{f(x)} = 
\lim\limits_{x\rightarrow0+}\dfrac{-(1/2)x}{\sqrt{x}} =
-\dfrac{1}{2}\cdot\lim\limits_{x\rightarrow0+}x^{\frac{1}{2}} = 0. \\
$$
$\therefore$\ $g(x)$ is infinitely-small relative to $f(x)$, as $x\rightarrow0+$. In Bachmann-Landau notation this may be represented as $g(x)=o(f(x))$.\\ [10pt]
Moreover, it is possible to calculate the exact asymptotic order of growth of one function against the other as $x\rightarrow0+$. Let us consider the following relation
$$
\lim\limits_{x\rightarrow0+}\dfrac{g(x)}{(f(x))^\alpha} = C,
$$
where $\alpha$ is the asymptotic order of growth of g(x) relative to f(x) and $C$ is a constant.
$$
\begin{array}{c}
\lim\limits_{x\rightarrow0+}\dfrac{g(x)}{(f(x))^\alpha} = C,\\ [14pt]
\lim\limits_{x\rightarrow0+}\dfrac{-(1/2)x}{(\sqrt{x})^\alpha} = C,\\ [12pt]
-\dfrac{1}{2}\cdot\lim\limits_{x\rightarrow0+}\dfrac{x}{x^{\frac{\alpha}{2}}} = C, \\ [10pt]
\therefore\ \alpha=2.
\end{array}
$$

% ---------------------------- Task 4 ----------------------------------
\newpage
\subsection*{\center Question \#4.}
{\bf Conditions:~}
Find the points of discontinuity of the function $y=f(x)$ and state their class. Plot the graph of $f(x)$ in the neighborhood of each point of discontinuity.
$$
f(x) \equiv 
\begin{cases}
\cos{\left(\dfrac{\pi x}{2}\right)},	& |x|\leqslant1, \\
|x-1|,   & |x|>1.
\end{cases}
$$
{\bf Solution:}
Both $\cos{\left(\dfrac{\pi x}{2}\right)}$ and $|x-1|$ are continuous functions on $\mathbb{R}$, therefore a discontinuity may only occur in the points $x_1=-1$ and $x_2=1$. As such, we should evaluate the one-sided limits of $f(x)$ in these two points:
$$
\begin{array}{c}
\lim\limits_{x\rightarrow(-1)-}f(x) = 
\lim\limits_{x\rightarrow(-1)-}|x-1| = 2, \qquad
\lim\limits_{x\rightarrow(-1)+}f(x) = 
\lim\limits_{x\rightarrow(-1)+}\cos{\left(\dfrac{\pi x}{2}\right)} = 0, \\[12pt]
\therefore\text{ Both limits exist and are finite, yet }\lim\limits_{x\rightarrow(-1)-}f(x)\neq \lim\limits_{x\rightarrow(-1)+}f(x), \\
\implies x_1=-1, \text{ is a point of jump discontinuity.} \\[16pt]
\lim\limits_{x\rightarrow1-}f(x) = 
\lim\limits_{x\rightarrow1-}\cos{\left(\dfrac{\pi x}{2}\right)} = 0, \qquad
\lim\limits_{x\rightarrow1+}f(x) = 
\lim\limits_{x\rightarrow1+}|x-1| = 0, \\
\therefore\text{ Both limits exist, are finite and are equal to 0,} \\
\implies x_2=1, \text{ is not a point of discontinuity.} \\[15pt]
\end{array}
$$
Knowing the points of discontinuity and their classification, it is possible to plot the graph of $f(x)$ in the neighborhood of the points of discontinuity:
\vspace{0.2cm}
\begin{center}
	\begin{tikzpicture}
	\def\outerfunc{}
	\begin{axis}[xmin=-5.0,
	xmax=5.0, 
	ymin=-0.8,
	ymax=4.5,
	width=\textwidth,
	height=0.61\textwidth,
	axis x line=middle,
	axis y line=middle, 
	every axis x label/.style={at={(current axis.right of origin)},anchor=west},
	every inner x axis line/.append style={|-latex'},
	every inner y axis line/.append style={|-latex'},
	minor tick num=1,			
	axis equal=true,
	xlabel=$x$, 
	ylabel=$y$,          
	samples=800,
	clip=true,
	]
	\addplot[color=black, line width=1.4pt,domain=-5:-1] {-\x+1};
	\addplot[color=black, line width=1.4pt,domain=-1:1]{cos(deg((\x*pi)/2))};
	\addplot[color=black, line width=1.4pt,domain=1:5]{\x-1};
	\addplot[
	mark=*,
	mark options={fill=black, draw=black},
	] coordinates {(-1, 0) (1, 0)};
	\addplot[
	mark=*,
	mark options={fill=white, draw=black},
	] coordinates {(-1, 2)};
	\end{axis}
	\end{tikzpicture}
\end{center}


\newpage
\addcontentsline{toc}{section}{Bibliography}
\begin{thebibliography}{3}
\bibitem{book01} L'vovsky S.M. Typesetting and layout in the \LaTeX\ system, 2003 c.
\bibitem{website1} \LaTeX\ Base Reference, https://docs.latexbase.com/, 2020.
\bibitem{website2} \LaTeXe\ Cheat Sheet, https://wch.github.io/latexsheet/latexsheet.pdf, 2014.
\end{thebibliography}

\end{document}
